% $RUN6 = paper/data/run8-randomization-window
\section{The Hypotheses Read Back}
\label{sec:verdicts}
Section~\ref{sec:hypotheses} states the six hypotheses and the three verdict words. This section says when the hypotheses were written, reads them back in Table~\ref{tab:verdicts}, and names what held and what did not.

\paragraph{When they were written.}
\begin{itemize}
% src: git log of paper/plan.md: one commit, 145eaf7 (3 November 2026), which also holds paper/data/run2-prelaunch-clean/gates.json and run-log.json; the plan's section 1a reads the first real run
\item Commit 145eaf7 (3 November 2026) holds the plan beside the reports of the second run on the same split, and the plan's section 1a reads the first. So the hypotheses postdate those two runs: the cautionary run of Section~\ref{sec:cautionary} and the pre-launch-clean run.
% src: references.bib mootcourt2026plan: tag plan-registered-2026-11-04, commit f0b088c794dcea1a1eac9872cd15aeb1c2412976, pushed 2026-11-04T17:19:12Z, plan.md sha256 e6dab9062fd1dc59734e5e4374ba429d3edc06efe681f68bb81424d88f45b613, byte-identical to paper/plan.md of 145eaf7
\item A byte-identical copy of the plan was posted publicly on 4 November 2026, after those runs and after the first review of a draft \cite{mootcourt2026plan}. The commit and the tag named in that reference are the record.
% src: paper/plan.md line 31 (H5: "more positive lift than random"), the plan's E4 and E12; docs/method.md call 201 (the arm-count-matched order); this paper's Section 5 (the reading rule)
\item Four readings postdate the data: the pooled interval as the primary claim, the arm-count-matched baseline of the fifth hypothesis, reading that hypothesis over the challengers as well as the ladder, and the verdict words themselves. On the ladder alone the fifth hypothesis does not hold under the matched baseline.
% src: docs/superpowers/plans/2026-11-06-paper-v2-plan.md D1, D2 (the randomization window, 6 November 2026)
\item The date rule of Section~\ref{sec:setting} changed the rows these hypotheses are read on after they were written, and it did not change their wording.
\end{itemize}
Written after two runs on the same split and read by rules partly fixed after the data, these hypotheses do not make this a pre-registered study, and the verdicts below are exploratory like every other number in the paper.

\begin{table*}[t]\centering\scriptsize
% src: $RUN6/backtest-upworthy-20261106-2dc461f52ea4.json and $RUN6/backtest-asos-20261103-6bb5a15c5ff0.json (the main backtests); $RUN6/paper-extras-upworthy-20261106-2dc461f52ea4.json keys p1_minus_p0_by_origin, issued_ece_by_origin, issued_row_brier_skill, abstention_composition and drift_matched_null, bootstrap_resamples 500; $RUN6/learning-curve-upworthy-20261106-2dc461f52ea4.json; $RUN6/paper-extras-synthetic-20261103-bf4fd286e2b2.json positive_control_rungs; tables.decision_value budget=0.25, family_size 6; tables/verdicts.tex
\caption{The six hypotheses read back part by part under the rule of Section~\ref{sec:hypotheses}. Rows: the parts of each hypothesis, under a line that gives its overall reading and the section that holds the result; H5 is one part, read rung by rung, and H2 is read two ways. Columns: what the plan fixed, in short (Section~\ref{sec:hypotheses} gives the wording), what was measured, and the verdict on the part. A pair of numbers after a value is its 90 percent bootstrap interval by test; an adjusted interval or bound is the Bonferroni-adjusted one of the decision table. H1 reads the log score in nats, H2 and H4 the Brier skill against the base rate, and H5 the captured lift over the baseline at a quarter of the tests.}
\label{tab:verdicts}
% generated by paper/build_tables.py from the report keys named above each row (limits)
\begin{tabular}{@{}p{0.10\textwidth}p{0.21\textwidth}p{0.45\textwidth}p{0.17\textwidth}@{}}
\toprule
\raggedright hypothesis & \raggedright what was planned & \raggedright what was measured & \raggedright verdict \tabularnewline
\midrule
\multicolumn{4}{@{}l}{\textbf{H1, in part} (see~\ref{sec:ladder-results})} \\
% src: $RUN6/backtest-upworthy-20261106-2dc461f52ea4.json tables.score_by_origin[origin=2015-01-01T00:00:00+00:00,rung=P1].logscore; $RUN6/backtest-upworthy-20261106-2dc461f52ea4.json tables.score_by_origin[origin=2015-01-01T00:00:00+00:00,rung=P0].logscore; $RUN6/paper-extras-upworthy-20261106-2dc461f52ea4.json p1_minus_p0_by_origin.by_origin[origin=2015-01-01T00:00:00+00:00].difference; $RUN6/paper-extras-upworthy-20261106-2dc461f52ea4.json p1_minus_p0_by_origin.by_origin[origin=2015-01-01T00:00:00+00:00].lo; $RUN6/paper-extras-upworthy-20261106-2dc461f52ea4.json p1_minus_p0_by_origin.by_origin[origin=2015-01-01T00:00:00+00:00].hi; $RUN6/paper-extras-upworthy-20261106-2dc461f52ea4.json p1_minus_p0_by_origin.by_origin_clear_of_zero
\raggedright H1 part 1 & \raggedright P1 beats P0 on log score at every origin after burn-in, P0 scored with a normal density, as in the code committed with the plan & \raggedright above at every origin after burn-in; the narrowest margin January 2015, minus 0.730 against minus 0.761; paired by test the difference there is plus 0.031, 0.007 to 0.055, and clear of zero at all seven origins & \raggedright held \tabularnewline
% src: $RUN6/backtest-upworthy-20261106-2dc461f52ea4.json tables.aggregate[cell=headline_only,rung=P2].logscore_vs_p1; $RUN6/backtest-upworthy-20261106-2dc461f52ea4.json tables.aggregate[cell=headline_only,rung=P2].logscore_vs_p1_lo; $RUN6/backtest-upworthy-20261106-2dc461f52ea4.json tables.aggregate[cell=headline_only,rung=P2].logscore_vs_p1_hi; $RUN6/paper-extras-synthetic-20261103-bf4fd286e2b2.json positive_control_rungs.runs[n_tests=11100].by_cell["headline_only"].rungs[rung=P2].diff_logscore_vs_p1
\raggedright H1 part 2 & \raggedright P2 beats P1 on headline\_only & \raggedright plus 0.0002, minus 0.0007 to 0.0010; the interval contains the 0.0004 that a planted lift returns in a cell of this size under the same backtest & \raggedright not shown \tabularnewline
% src: $RUN6/backtest-upworthy-20261106-2dc461f52ea4.json tables.aggregate[cell=image_only,rung=P2].logscore_vs_p1; $RUN6/backtest-upworthy-20261106-2dc461f52ea4.json tables.aggregate[cell=image_only,rung=P2].logscore_vs_p1_lo; $RUN6/backtest-upworthy-20261106-2dc461f52ea4.json tables.aggregate[cell=image_only,rung=P2].logscore_vs_p1_hi
\raggedright H1 part 3 & \raggedright P2 does not beat P1 on image\_only & \raggedright plus 0.0000 nats, 0.0000 to 0.0000: the guard admitted no weight in this cell at any origin & \raggedright held \tabularnewline
\midrule
\multicolumn{4}{@{}l}{\textbf{H2, in part: the gate claim failed; skill not shown read pooled and undecidable as written; above zero at one earlier origin, where the base rate lags}} \\
% src: $RUN6/backtest-upworthy-20261106-2dc461f52ea4.json tables.aggregate[cell=ALL,rung=P2].brier_skill; $RUN6/backtest-upworthy-20261106-2dc461f52ea4.json tables.aggregate[cell=ALL,rung=P2].brier_skill_lo; $RUN6/backtest-upworthy-20261106-2dc461f52ea4.json tables.aggregate[cell=ALL,rung=P2].brier_skill_hi
\raggedright H2, read pooled & \raggedright P2's Brier skill positive, its interval clear of zero, pooled over origins & \raggedright 0.0032, minus 0.0004 to 0.0072, not clear of zero & \raggedright not shown \tabularnewline
% src: $RUN6/learning-curve-upworthy-20261106-2dc461f52ea4.json rows[fraction=1.0,rung=P2].skill; $RUN6/learning-curve-upworthy-20261106-2dc461f52ea4.json rows[fraction=1.0,rung=P2].skill_lo; $RUN6/learning-curve-upworthy-20261106-2dc461f52ea4.json rows[fraction=1.0,rung=P2].skill_hi; $RUN6/learning-curve-upworthy-20261106-2dc461f52ea4.json heldout_tests; $RUN6/paper-extras-upworthy-20261106-2dc461f52ea4.json gates_by_stratum.strata.ALL.own_gate.n_400; $RUN6/paper-extras-upworthy-20261106-2dc461f52ea4.json drift_matched_null.rungs.P0.by_origin["2015-01-01T00:00:00+00:00"].mean_forecast; $RUN6/paper-extras-upworthy-20261106-2dc461f52ea4.json drift_matched_null.rungs.P0.by_origin["2015-01-01T00:00:00+00:00"].event_rate; $RUN6/backtest-upworthy-20261106-2dc461f52ea4.json tables.origins[origin=2015-04-01T00:00:00+00:00].n_tests
\raggedright H2, as written & \raggedright the same at the final origin & \raggedright at January 2015, the latest origin with enough held-out rows for the learning curve: 0.029, 0.014 to 0.051, on 318 tests, where the pooled gate of 300 was simulated at the pooled event rate and not at this origin's; there the base rate is 0.091 and the event rate 0.058, so the base rate lags and part of the skill may be that lag; at April 2015, the last origin of the calendar: 83 tests, below every gate & \raggedright above zero at January 2015, with that caution, not counted as holding; undecidable at April 2015 \tabularnewline
% src: $RUN6/backtest-upworthy-20261106-2dc461f52ea4.json tables.origins[origin=2015-04-01T00:00:00+00:00].n_tests; $RUN6/paper-extras-upworthy-20261106-2dc461f52ea4.json gates_by_stratum.strata.ALL.own_gate.n_400
\raggedright H2, the gate & \raggedright the released gate decidable at the final origin & \raggedright the last origin of the calendar, April 2015, holds 83 tests, below the gate of 300 for a skill of 0.10 and below every other gate & \raggedright failed \tabularnewline
\midrule
\multicolumn{4}{@{}l}{\textbf{H3, in part} (see~\ref{sec:calibration-results})} \\
% src: $RUN6/paper-extras-upworthy-20261106-2dc461f52ea4.json issued_ece_by_origin.issued.origins_inside_central_90; $RUN6/paper-extras-upworthy-20261106-2dc461f52ea4.json issued_ece_by_origin.issued.origins; $RUN6/paper-extras-upworthy-20261106-2dc461f52ea4.json issued_ece_by_origin.issued.pooled.ece_percentile
\raggedright H3 part 1 & \raggedright ECE of issued cards inside the central 90 percent of the null at most origins & \raggedright inside at 2 of the 5 origins that issue cards, not at most of them; pooled at the 100th percentile of the clustered null & \raggedright failed \tabularnewline
% src: $RUN6/backtest-upworthy-20261106-2dc461f52ea4.json tables.aggregate[cell=ALL,rung=P1].coverage_80; $RUN6/backtest-upworthy-20261106-2dc461f52ea4.json tables.aggregate[cell=ALL,rung=P1].coverage_80_lo; $RUN6/backtest-upworthy-20261106-2dc461f52ea4.json tables.aggregate[cell=ALL,rung=P1].coverage_80_hi; $RUN6/backtest-upworthy-20261106-2dc461f52ea4.json tables.aggregate[cell=ALL,rung=P1].coverage_90; $RUN6/backtest-upworthy-20261106-2dc461f52ea4.json tables.aggregate[cell=ALL,rung=P1].coverage_90_lo; $RUN6/backtest-upworthy-20261106-2dc461f52ea4.json tables.aggregate[cell=ALL,rung=P1].coverage_90_hi; $RUN6/backtest-upworthy-20261106-2dc461f52ea4.json tables.coverage_by_se_ratio[rung=P1,decile=1].coverage_90; $RUN6/backtest-upworthy-20261106-2dc461f52ea4.json tables.coverage_by_se_ratio[rung=P1,decile=10].coverage_90
\raggedright H3 part 2 & \raggedright 80 and 90 percent coverage within 5 points of nominal & \raggedright 81 percent, 80 to 82, and 90 percent, 89 to 90, within 5 points of nominal pooled; 1.00 to 0.69 across deciles (a range, not an interval) of the ratio of realized to planned standard error & \raggedright held, read pooled; the fall across deciles is descriptive \tabularnewline
\midrule
\multicolumn{4}{@{}l}{\textbf{H4, in part} (see~\ref{sec:abstention-results})} \\
% src: $RUN6/paper-extras-upworthy-20261106-2dc461f52ea4.json issued_row_brier_skill.rungs.P2.issued.brier_skill; $RUN6/paper-extras-upworthy-20261106-2dc461f52ea4.json issued_row_brier_skill.rungs.P2.issued.brier_skill_lo; $RUN6/paper-extras-upworthy-20261106-2dc461f52ea4.json issued_row_brier_skill.rungs.P2.issued.brier_skill_hi; $RUN6/paper-extras-upworthy-20261106-2dc461f52ea4.json issued_row_brier_skill.rungs.P2.all.brier_skill
\raggedright H4 part 1 & \raggedright issued cards above the full set on skill & \raggedright issued rows 0.0004, minus 0.0035 to 0.0046, against 0.0032 for the full set & \raggedright failed \tabularnewline
% src: $RUN6/paper-extras-upworthy-20261106-2dc461f52ea4.json issued_row_brier_skill.rungs.P2.abstained.brier_skill; $RUN6/paper-extras-upworthy-20261106-2dc461f52ea4.json issued_row_brier_skill.rungs.P2.abstained.brier_skill_lo; $RUN6/paper-extras-upworthy-20261106-2dc461f52ea4.json issued_row_brier_skill.rungs.P2.abstained.brier_skill_hi; $RUN6/paper-extras-upworthy-20261106-2dc461f52ea4.json abstention_composition.P2.abstained.skill; $RUN6/paper-extras-upworthy-20261106-2dc461f52ea4.json abstention_composition.P2.abstained.lo; $RUN6/paper-extras-upworthy-20261106-2dc461f52ea4.json abstention_composition.P2.abstained.hi; $RUN6/paper-extras-upworthy-20261106-2dc461f52ea4.json abstention_composition.P2.issued.skill; $RUN6/paper-extras-upworthy-20261106-2dc461f52ea4.json abstention_composition.P2.issued.lo; $RUN6/paper-extras-upworthy-20261106-2dc461f52ea4.json abstention_composition.P2.issued.hi; $RUN6/paper-extras-upworthy-20261106-2dc461f52ea4.json abstention_composition.P2.pooled_issued_minus_abstained_skill; $RUN6/paper-extras-upworthy-20261106-2dc461f52ea4.json abstention_composition.P2.stratified_issued_minus_abstained_skill.estimate; $RUN6/paper-extras-upworthy-20261106-2dc461f52ea4.json abstention_composition.P2.stratified_issued_minus_abstained_skill.lo; $RUN6/paper-extras-upworthy-20261106-2dc461f52ea4.json abstention_composition.P2.stratified_issued_minus_abstained_skill.hi
\raggedright H4 part 2 & \raggedright abstained rows lower on skill when forced & \raggedright abstained rows, scored on the withheld forecast: with the burn-in origin, above the full set at 0.008, 0.001 to 0.015; without it, minus 0.0013, minus 0.0060 to 0.0041, against 0.0004, minus 0.0035 to 0.0045 for issued rows, no measurable difference & \raggedright failed on all rows; with burn-in excluded, on the claim's side pooled with no interval and on the wrong side within origin and cell \tabularnewline
% src: $RUN6/backtest-upworthy-20261106-2dc461f52ea4.json tables.abstention[cell=ALL].rate; $RUN6/backtest-upworthy-20261106-2dc461f52ea4.json tables.abstention[cell=ALL].rate; $RUN6/backtest-upworthy-20261106-2dc461f52ea4.json tables.abstention[cell=ALL].n; $RUN6/backtest-upworthy-20261106-2dc461f52ea4.json tables.origins[origin=2013-04-01T00:00:00+00:00].n_scored; $RUN6/backtest-upworthy-20261106-2dc461f52ea4.json tables.origins[origin=2013-04-01T00:00:00+00:00].abstention_rate
\raggedright H4 part 3 & \raggedright abstention below 30 percent after burn-in & \raggedright 32 percent pooled and 27.2 percent after the burn-in origin, January 2014 included & \raggedright held \tabularnewline
\midrule
\multicolumn{4}{@{}l}{\textbf{H5, in part} (Sections~\ref{sec:hypotheses} and~\ref{sec:decision-results})} \\
% src: $RUN6/backtest-upworthy-20261106-2dc461f52ea4.json tables.decision_value[origin=ALL,budget=0.25,policy=p_win,rung=text_lr].diff_matched; $RUN6/backtest-upworthy-20261106-2dc461f52ea4.json tables.decision_value[origin=ALL,budget=0.25,policy=p_win,rung=text_lr].diff_matched_adjusted_lo; $RUN6/backtest-upworthy-20261106-2dc461f52ea4.json tables.decision_value[origin=ALL,budget=0.25,policy=p_win,rung=text_lr].diff_matched_adjusted_hi
\raggedright H5, text classifier & \raggedright more positive lift than random at K = 25 percent of tests & \raggedright against the arm-count-matched random order, plus 0.046, adjusted interval above zero & \raggedright held against the matched order \tabularnewline
% src: $RUN6/backtest-upworthy-20261106-2dc461f52ea4.json tables.decision_value[origin=ALL,budget=0.25,policy=p_win,rung=GBM].diff_matched; $RUN6/backtest-upworthy-20261106-2dc461f52ea4.json tables.decision_value[origin=ALL,budget=0.25,policy=p_win,rung=GBM].diff_matched_adjusted_lo; $RUN6/backtest-upworthy-20261106-2dc461f52ea4.json tables.decision_value[origin=ALL,budget=0.25,policy=p_win,rung=GBM].diff_matched_adjusted_hi; $RUN6/backtest-upworthy-20261106-2dc461f52ea4.json tables.decision_value[origin=ALL,budget=0.25,policy=p_win,rung=GBM].diff_random_adjusted_lo; $RUN6/backtest-upworthy-20261106-2dc461f52ea4.json tables.decision_value[origin=ALL,budget=0.25,policy=p_win,rung=GBM].diff_random_adjusted_hi
\raggedright H5, boosted model & \raggedright the same & \raggedright against the matched order, plus 0.017, adjusted interval spanning zero; against the random order the plan names, adjusted 0.001 to 0.052 & \raggedright not shown against the matched order; held against the plan's random order \tabularnewline
% src: $RUN6/backtest-upworthy-20261106-2dc461f52ea4.json tables.decision_value[origin=ALL,budget=0.25,policy=p_win,rung=P2].diff_matched; $RUN6/backtest-upworthy-20261106-2dc461f52ea4.json tables.decision_value[origin=ALL,budget=0.25,policy=p_win,rung=P2].diff_matched_adjusted_lo; $RUN6/backtest-upworthy-20261106-2dc461f52ea4.json tables.decision_value[origin=ALL,budget=0.25,policy=p_win,rung=P2].diff_matched_adjusted_hi; $RUN6/backtest-upworthy-20261106-2dc461f52ea4.json tables.decision_value[origin=ALL,budget=0.25,policy=p_win,rung=P2].diff_random_adjusted_lo; $RUN6/backtest-upworthy-20261106-2dc461f52ea4.json tables.decision_value[origin=ALL,budget=0.25,policy=p_win,rung=P2].diff_random_adjusted_hi
\raggedright H5, precedent rung & \raggedright the same & \raggedright against the matched order, plus 0.018, adjusted interval spanning zero; against the random order the plan names, adjusted 0.004 to 0.045 & \raggedright not shown against the matched order; held against the plan's random order \tabularnewline
% src: $RUN6/backtest-upworthy-20261106-2dc461f52ea4.json tables.decision_value[origin=ALL,budget=0.25,policy=p_win,rung=local_eb].diff_matched; $RUN6/backtest-upworthy-20261106-2dc461f52ea4.json tables.decision_value[origin=ALL,budget=0.25,policy=p_win,rung=local_eb].diff_matched_adjusted_lo; $RUN6/backtest-upworthy-20261106-2dc461f52ea4.json tables.decision_value[origin=ALL,budget=0.25,policy=p_win,rung=local_eb].diff_matched_adjusted_hi; $RUN6/backtest-upworthy-20261106-2dc461f52ea4.json tables.decision_value[origin=ALL,budget=0.25,policy=p_win,rung=local_eb].diff_random_adjusted_lo; $RUN6/backtest-upworthy-20261106-2dc461f52ea4.json tables.decision_value[origin=ALL,budget=0.25,policy=p_win,rung=local_eb].diff_random_adjusted_hi
\raggedright H5, local comparison set & \raggedright the same & \raggedright against the matched order, plus 0.018, adjusted interval spanning zero; against the random order the plan names, adjusted 0.002 to 0.056 & \raggedright not shown against the matched order; held against the plan's random order \tabularnewline
% src: $RUN6/backtest-upworthy-20261106-2dc461f52ea4.json tables.decision_value[origin=ALL,budget=0.25,policy=p_win,rung=P1].diff_matched; $RUN6/backtest-upworthy-20261106-2dc461f52ea4.json tables.decision_value[origin=ALL,budget=0.25,policy=p_win,rung=P1].diff_matched_lo; $RUN6/backtest-upworthy-20261106-2dc461f52ea4.json tables.decision_value[origin=ALL,budget=0.25,policy=p_win,rung=P1].diff_matched_hi
\raggedright H5, prior & \raggedright the same & \raggedright at the matched order within noise, minus 0.004, minus 0.014 to 0.007 & \raggedright failed \tabularnewline
\midrule
\multicolumn{4}{@{}l}{\textbf{H6, held} (see~\ref{sec:asos-results})} \\
% src: $RUN6/backtest-asos-20261103-6bb5a15c5ff0.json tables.aggregate[cell=ALL,rung=P2].n_tests
\raggedright H6 part 1 & \raggedright every gate undecidable on the retailer's archive & \raggedright undecidable at every gate & \raggedright held, a result 68 tests could not have avoided \tabularnewline
\raggedright H6 part 2 & \raggedright P2 equals P1 & \raggedright P2 equals P1 on this archive, because the guard holds the weight at zero & \raggedright held \tabularnewline
\bottomrule
\end{tabular}

\end{table*}

\paragraph{What held and what did not.} The prior beat the base rate on the log score at every origin after burn-in, abstention stayed under its limit after burn-in, the text classifier captured more lift than the matched order, and the retailer's archive stayed undecidable. Pooled skill at realized precision was not shown, calibration of issued cards failed, issued cards did not score above the full set on Brier skill, and the prior did not rank tests. A precedent gain in the largest cell was not shown, and the backtest could not have shown one of the size a planted control returns. Three parts hold for reasons that limit what they say, as the notes below explain.

\paragraph{Notes on the table.}
\begin{itemize}
\item H1 part 2: the interval contains what a planted signal of that size returns, so this result does not refute a precedent signal; it says this backtest cannot see one (Section~\ref{sec:synthetic-results}).
\item H1 part 3: P2 equals P1 in the image-only cell because the guard admitted no weight there, so the part holds by the guard and says nothing about what precedents carry in that cell.
% src: $RUN6/learning-curve-upworthy-20261106-2dc461f52ea4.json rows[fraction=1.0,rung=P2] skill 0.029 (0.014, 0.051); $RUN6/paper-extras-upworthy-20261106-2dc461f52ea4.json drift_matched_null.rungs.P0.by_origin["2015-01-01T00:00:00+00:00"] mean_forecast 0.091, event_rate 0.058, drift_matched_null.rungs.P2.by_origin["2015-01-01T00:00:00+00:00"] mean_forecast 0.066; within_origin_auc_by_year.years["2015"].P2.auc_within_origin 0.634 (0.568, 0.703)
\item H2: the plan says ``at the final origin'', and the verdict depends on the reading. Pooled, under the rule of Section~\ref{sec:hypotheses}, it is not shown.
% src: $RUN6/backtest-upworthy-20261106-2dc461f52ea4.json tables.aggregate[cell=ALL,rung=P2] brier_skill_planned 0.0058 (0.0018, 0.0095)
With the event at planned precision instead, P2's pooled interval lies above zero, unadjusted, at 0.0058, so this reading depends on how the event is defined and the skill is at most small either way (Section~\ref{sec:ladder-results}).
As written it holds at January 2015, the last origin with enough tests to clear a gate, and the last calendar origin is undecidable. At January 2015 the base rate of 0.091 stands above an event rate of 0.058 and P2's mean forecast of 0.066 sits near the event rate, so part of that skill may be the base rate's lag rather than foresight about any one test (Table~\ref{tab:fmt-auc-by-year} gives the ranking within that year).
\item H3: the ECE percentile rejects calibration, and Section~\ref{sec:calibration-results} shows that most of the error is a level the forecast does not anticipate.
\item H4: the hypothesis names skill. With the burn-in origin the abstained rows score above the full set on it, and without it the abstained and issued rows do not differ measurably (Table~\ref{tab:fmt-abstention-composition}). On the log score, which no part names, issued rows are better (Section~\ref{sec:abstention-results}). The 27.2 percent after burn-in includes January 2014, which abstains on every row because the date rule leaves its reliability reports stale.
\item H5: the prior ranks by cell only, its win probability one number per cell and origin. The verdict rests on the arm-count-matched order, adopted after the data, and regret against the oracle is reported as is in Section~\ref{sec:decision-results}.
\item H6: 68 tests against gates of hundreds could not have come out otherwise, and P2 equals P1 only because the guard holds the weight at zero. This hypothesis confirms the protocol's behaviour on a small archive and is no evidence about product experiments.
\end{itemize}
